% This is text associated with the open textbook "Open Electromagnetics"
% See immediately below for author, date
% This work is licensed under the Creative Commons Attribution-ShareAlike 4.0 International License. (CC BY-SA 4.0) 
% To view a copy of the license, visit http://creativecommons.org/licenses/by-sa/4.0/

%ORB TITLE Parallel Plate Waveguide: TE Case, Magnetic Field
%ORB AUTHOR 0        # S.W. Ellingson, Virginia Tech (ellingson@vt.edu)
%ORB DATE 20191121   # yyyymmdd
%ORB ABSTRACT 

\label{m0175_Parallel_Plate_Waveguide_TE_Magnetic} % <-- Must match name of this file and module directory name.  Don't mess with this!
\index{waveguide!parallel plate|(}
\index{parallel-plate waveguide|see{waveguide}}

In Section~\ref{m0173_Parallel_Plate_Waveguide_Introduction}, the parallel plate waveguide was introduced.
In Section~\ref{m0174_Parallel_Plate_Waveguide_TE_Electric}, we determined the TE component of the electric field.
In this section, we determine the TE component of the magnetic field.
The reader should be familiar with Section~\ref{m0174_Parallel_Plate_Waveguide_TE_Electric} before attempting this section.

In Section~\ref{m0174_Parallel_Plate_Waveguide_TE_Electric}, the TE component of the electric field was determined to be:
%
\begin{equation}
\hat{\bf y}\widetilde{E}_y = \hat{\bf y}\sum_{m=1}^{\infty} \widetilde{E}_y^{(m)}
\label{m0175_eEysum}
\end{equation}
%
The TE component of the magnetic field may be obtained from the Maxwell-Faraday equation:
\index{Maxwell-Faraday equation}
%
\begin{equation}
\nabla \times \widetilde{\bf E} = -j\omega\mu \widetilde{\bf H}
\end{equation}
%
Thus:
%
\begin{flalign}
\widetilde{\bf H} &= \frac{j}{\omega\mu}~\nabla \times \widetilde{\bf E} \nonumber \\
                  &= \frac{j}{\omega\mu}~\nabla \times \left(\hat{\bf y}\widetilde{E}_y\right) 
\end{flalign}
%
The relevant form of the curl operator is Equation~\ref{m0139_eCurlCart} (Appendix~\ref{m0139_Vector_Operators}).
Although the complete expression consists of 6 terms, all but 2 of these terms are zero because the $\hat{\bf x}$ and $\hat{\bf z}$ components of $\widetilde{\bf E}$ are zero.
The two remaining terms are 
$-\hat{\bf x}\partial \widetilde{E}_y/\partial z$ and
$+\hat{\bf z}\partial \widetilde{E}_y/\partial x$.
Thus:
%
\begin{equation}
\widetilde{\bf H} = \frac{j}{\omega\mu}~\left( -\hat{\bf x}\frac{\partial\widetilde{E}_y}{\partial z} +\hat{\bf z}\frac{\partial\widetilde{E}_y}{\partial x} \right) 
\label{m0175_eH}
\end{equation}
%
Recall that $\widetilde{E}_y$ is the sum of modes, as indicated in Equation~\ref{m0175_eEysum}.
Since differentiation (i.e., $\partial/\partial z$ and $\partial/\partial x$) is a linear operator, we may evaluate Equation~\ref{m0175_eH} for modes one at a time, and then sum the results.
Using this approach, we find:
%
\begin{flalign}
\frac{\partial\widetilde{E}_y^{(m)}}{\partial z} &= \frac{\partial}{\partial z} ~ E_{y0}^{(m)} e^{-jk_z^{(m)} z} \sin k_x^{(m)} x \nonumber \\
                                         &= \left( E_{y0}^{(m)} e^{-jk_z^{(m)} z} \sin k_x^{(m)} x \right) \left(-jk_z^{(m)}\right) 
\end{flalign}
%
and
%
\begin{flalign}
\frac{\partial\widetilde{E}_y^{(m)}}{\partial x} &= \frac{\partial}{\partial x} ~ E_{y0}^{(m)} e^{-jk_z^{(m)} z} \sin k_x^{(m)} x \nonumber \\
                                         &= \left( E_{y0}^{(m)} e^{-jk_z^{(m)} z} \cos k_x^{(m)} x \right) \left(+k_x^{(m)}\right) 
\end{flalign}
%
We may now assemble a solution for the magnetic field as follows:
%
\begin{flalign}
\hat{\bf x}\widetilde{H}_x+\hat{\bf z}\widetilde{H}_z =&~~~\hat{\bf x}\sum_{m=1}^{\infty} \widetilde{H}_x^{(m)} \nonumber \\
                                                       &+\hat{\bf z}\sum_{m=1}^{\infty} \widetilde{H}_z^{(m)} \label{m0175_eHsum}
\end{flalign}
%
where
%
\begin{equation}
\widetilde{H}_x^{(m)} = -\frac{k_z^{(m)}}{\omega\mu} E_{y0}^{(m)} e^{-jk_z^{(m)} z} \sin k_x^{(m)} x 
\label{m0175_eHx}
\end{equation}
%
\begin{equation}
\widetilde{H}_z^{(m)} = +j\frac{k_x^{(m)}}{\omega\mu} E_{y0}^{(m)} e^{-jk_z^{(m)} z} \cos k_x^{(m)} x
\label{m0175_eHz} 
\end{equation}
%
and modes may only exist at frequencies greater than the associated cutoff frequencies.  Summarizing:
%
%%%%%%%%%%%%%%%%%%%%%%%%%%%%%%%%%%%%%%%%%%%%%%%
%ORB HIGHLIGHT
\begin{mdframed}[backgroundcolor=yellow!20,nobreak=true] 
The magnetic field component of the TE solution is given by Equation~\ref{m0175_eHsum} with modal components as indicated by Equations~\ref{m0175_eHx} and \ref{m0175_eHz}.  Caveats pertaining to the cutoff frequencies and the locations of sources and structures continue to apply.
\end{mdframed}
%%%%%%%%%%%%%%%%%%%%%%%%%%%%%%%%%%%%%%%%%%%%%%%

This result is quite complex, yet some additional insight is possible.
At the perfectly-conducting (PEC) surface at $x=0$, we see 
%
\begin{flalign}
\widetilde{H}_x^{(m)}(x=0) &= 0 \nonumber \\
\widetilde{H}_z^{(m)}(x=0) &= +j\frac{k_x^{(m)}}{\omega\mu} E_{y0}^{(m)} e^{-jk_z^{(m)} z} 
\end{flalign}
%
Similarly, on the PEC surface at $x=a$, we see 
%
\begin{flalign}
\widetilde{H}_x^{(m)}(x=a) &= 0 \nonumber \\
\widetilde{H}_z^{(m)}(x=a) &= -j\frac{k_x^{(m)}}{\omega\mu} E_{y0}^{(m)} e^{-jk_z^{(m)} z} 
\end{flalign}
%
Thus, we see the magnetic field vector at the PEC surfaces is non-zero and parallel to the PEC surfaces.
Recall that the magnetic field is identically zero \emph{inside} a PEC material.
Also recall that boundary conditions require that discontinuity in the component of ${\bf H}$ tangent to a surface must be supported by a surface current.
We conclude that
%
%%%%%%%%%%%%%%%%%%%%%%%%%%%%%%%%%%%%%%%%%%%%%%%
%ORB HIGHLIGHT
\begin{mdframed}[backgroundcolor=yellow!20,nobreak=true] 
Current flows on the PEC surfaces of the waveguide.
\end{mdframed}
%%%%%%%%%%%%%%%%%%%%%%%%%%%%%%%%%%%%%%%%%%%%%%%
%
If this seems surprising, note that essentially the same thing happens in a coaxial transmission line.
\index{transmission line!coaxial}
That is, signals in a coaxial transmission line can be described equally well in terms of either potentials and currents on the inner and outer conductors, or the electromagnetic fields between the conductors.
The parallel plate waveguide is only slightly more complicated because the field in a properly-designed coaxial cable is a single transverse electromagnetic (TEM) 
\index{transverse electromagnetic (TEM)}
mode, whereas the fields in a parallel plate waveguide are combinations of TE and TM modes.

Interestingly, we also find that the magnetic field vector points in different directions depending on position relative to the conducting surfaces. 
We just determined that the magnetic field is parallel to the conducting surfaces at those surfaces. 
However, the magnetic field is \emph{perpendicular} to those surfaces at $m$ locations between $x=0$ and $x=a$.
These locations correspond to maxima in the electric field.

\index{waveguide!parallel plate|)}


